\documentclass[a4paper,12pt,twoside]{memoir}

% Castellano
\usepackage[spanish,es-tabla]{babel}
\selectlanguage{spanish}
\usepackage[utf8]{inputenc}
\usepackage[T1]{fontenc}
\usepackage{lmodern} % scalable font
\usepackage{microtype}
\usepackage{placeins}

\RequirePackage{booktabs}
\RequirePackage[table]{xcolor}
\RequirePackage{xtab}
\RequirePackage{multirow}
% Bibliography management
\usepackage[numbers,sort]{natbib}


\usepackage{longtable}

% Links
\PassOptionsToPackage{hyphens}{url}\usepackage[colorlinks]{hyperref}
\hypersetup{
	allcolors = {red}
}

% Ecuaciones
\usepackage{amsmath}

% Rutas de fichero / paquete
\newcommand{\ruta}[1]{{\sffamily #1}}

% Párrafos
\nonzeroparskip

% Huérfanas y viudas
\widowpenalty100000
\clubpenalty100000

% Evitar solapes en el header
\nouppercaseheads

% Imagenes
\usepackage{graphicx}
\newcommand{\imagen}[2]{
	\begin{figure}[!h]
		\centering
		\includegraphics[width=0.9\textwidth]{#1}
		\caption{#2}\label{fig:#1}
	\end{figure}
	\FloatBarrier
}

\newcommand{\imagenflotante}[2]{
	\begin{figure}%[!h]
		\centering
		\includegraphics[width=0.9\textwidth]{#1}
		\caption{#2}\label{fig:#1}
	\end{figure}
}

\newcommand{\imagenConTamano}[3]{
	\begin{figure}[!h]
		\centering
		\includegraphics[width=#3\textwidth]{#1}
		\caption{#2}\label{fig:#1}
	\end{figure}
	\FloatBarrier
}

% El comando \figura nos permite insertar figuras comodamente, y utilizando
% siempre el mismo formato. Los parametros son:
% 1 -> Porcentaje del ancho de página que ocupará la figura (de 0 a 1)
% 2 --> Fichero de la imagen
% 3 --> Texto a pie de imagen
% 4 --> Etiqueta (label) para referencias
% 5 --> Opciones que queramos pasarle al \includegraphics
% 6 --> Opciones de posicionamiento a pasarle a \begin{figure}
\newcommand{\figuraConPosicion}[6]{%
  \setlength{\anchoFloat}{#1\textwidth}%
  \addtolength{\anchoFloat}{-4\fboxsep}%
  \setlength{\anchoFigura}{\anchoFloat}%
  \begin{figure}[#6]
    \begin{center}%
      \Ovalbox{%
        \begin{minipage}{\anchoFloat}%
          \begin{center}%
            \includegraphics[width=\anchoFigura,#5]{#2}%
            \caption{#3}%
            \label{#4}%
          \end{center}%
        \end{minipage}
      }%
    \end{center}%
  \end{figure}%
}

%
% Comando para incluir imágenes en formato apaisado (sin marco).
\newcommand{\figuraApaisadaSinMarco}[5]{%
  \begin{figure}%
    \begin{center}%
    \includegraphics[angle=90,height=#1\textheight,#5]{#2}%
    \caption{#3}%
    \label{#4}%
    \end{center}%
  \end{figure}%
}
% Para las tablas
\newcommand{\otoprule}{\midrule [\heavyrulewidth]}
%
% Nuevo comando para tablas pequeñas (menos de una página).
\newcommand{\tablaSmall}[5]{%
 \begin{table}
  \begin{center}
   \rowcolors {2}{gray!35}{}
   \begin{tabular}{#2}
    \toprule
    #4
    \otoprule
    #5
    \bottomrule
   \end{tabular}
   \caption{#1}
   \label{tabla:#3}
  \end{center}
 \end{table}
}

%
%Para el float H de tablaSmallSinColores
\usepackage{float}

%
% Nuevo comando para tablas pequeñas (menos de una página).
\newcommand{\tablaSmallSinColores}[5]{%
 \begin{table}[H]
  \begin{center}
   \begin{tabular}{#2}
    \toprule
    #4
    \otoprule
    #5
    \bottomrule
   \end{tabular}
   \caption{#1}
   \label{tabla:#3}
  \end{center}
 \end{table}
}

\newcommand{\tablaApaisadaSmall}[5]{%
\begin{landscape}
  \begin{table}
   \begin{center}
    \rowcolors {2}{gray!35}{}
    \begin{tabular}{#2}
     \toprule
     #4
     \otoprule
     #5
     \bottomrule
    \end{tabular}
    \caption{#1}
    \label{tabla:#3}
   \end{center}
  \end{table}
\end{landscape}
}

%
% Nuevo comando para tablas grandes con cabecera y filas alternas coloreadas en gris.
\newcommand{\tabla}[6]{%
  \begin{center}
    \tablefirsthead{
      \toprule
      #5
      \otoprule
    }
    \tablehead{
      \multicolumn{#3}{l}{\small\sl continúa desde la página anterior}\\
      \toprule
      #5
      \otoprule
    }
    \tabletail{
      \hline
      \multicolumn{#3}{r}{\small\sl continúa en la página siguiente}\\
    }
    \tablelasttail{
      \hline
    }
    \bottomcaption{#1}
    \rowcolors {2}{gray!35}{}
    \begin{xtabular}{#2}
      #6
      \bottomrule
    \end{xtabular}
    \label{tabla:#4}
  \end{center}
}

%
% Nuevo comando para tablas grandes con cabecera.
\newcommand{\tablaSinColores}[6]{%
  \begin{center}
    \tablefirsthead{
      \toprule
      #5
      \otoprule
    }
    \tablehead{
      \multicolumn{#3}{l}{\small\sl continúa desde la página anterior}\\
      \toprule
      #5
      \otoprule
    }
    \tabletail{
      \hline
      \multicolumn{#3}{r}{\small\sl continúa en la página siguiente}\\
    }
    \tablelasttail{
      \hline
    }
    \bottomcaption{#1}
    \begin{xtabular}{#2}
      #6
      \bottomrule
    \end{xtabular}
    \label{tabla:#4}
  \end{center}
}

%
% Nuevo comando para tablas grandes sin cabecera.
\newcommand{\tablaSinCabecera}[5]{%
  \begin{center}
    \tablefirsthead{
      \toprule
    }
    \tablehead{
      \multicolumn{#3}{l}{\small\sl continúa desde la página anterior}\\
      \hline
    }
    \tabletail{
      \hline
      \multicolumn{#3}{r}{\small\sl continúa en la página siguiente}\\
    }
    \tablelasttail{
      \hline
    }
    \bottomcaption{#1}
  \begin{xtabular}{#2}
    #5
   \bottomrule
  \end{xtabular}
  \label{tabla:#4}
  \end{center}
}



\definecolor{cgoLight}{HTML}{EEEEEE}
\definecolor{cgoExtralight}{HTML}{FFFFFF}

%
% Nuevo comando para tablas grandes sin cabecera.
\newcommand{\tablaSinCabeceraConBandas}[5]{%
  \begin{center}
    \tablefirsthead{
      \toprule
    }
    \tablehead{
      \multicolumn{#3}{l}{\small\sl continúa desde la página anterior}\\
      \hline
    }
    \tabletail{
      \hline
      \multicolumn{#3}{r}{\small\sl continúa en la página siguiente}\\
    }
    \tablelasttail{
      \hline
    }
    \bottomcaption{#1}
    \rowcolors[]{1}{cgoExtralight}{cgoLight}

  \begin{xtabular}{#2}
    #5
   \bottomrule
  \end{xtabular}
  \label{tabla:#4}
  \end{center}
}




\graphicspath{ {./img/} }

% Capítulos
\chapterstyle{bianchi}
\newcommand{\capitulo}[2]{
	\setcounter{chapter}{#1}
	\setcounter{section}{0}
	\chapter*{#2}
	\addcontentsline{toc}{chapter}{#2}
	\markboth{#2}{#2}
}

% Apéndices
\renewcommand{\appendixname}{Apéndice}
\renewcommand*\cftappendixname{\appendixname}

\newcommand{\apendice}[1]{
	%\renewcommand{\thechapter}{A}
	\chapter{#1}
}

\renewcommand*\cftappendixname{\appendixname\ }

\newcommand{\hubu}{\hline \rule{0pt}{2.7ex}}
% Formato de portada
\makeatletter
\usepackage{xcolor}
\newcommand{\tutor}[1]{\def\@tutor{#1}}
\newcommand{\course}[1]{\def\@course{#1}}
\definecolor{cpardoBox}{HTML}{E6E6FF}
\def\maketitle{
  \null
  \thispagestyle{empty}
  % Cabecera ----------------
\noindent\includegraphics[width=\textwidth]{cabecera}\vspace{1cm}%
  \vfill
  % Título proyecto y escudo informática ----------------
  \colorbox{cpardoBox}{%
    \begin{minipage}{.8\textwidth}
      \vspace{.5cm}\Large
      \begin{center}
      \textbf{TFG del Grado en Ingeniería Informática}\vspace{.6cm}\\
      \textbf{\LARGE\@title{}}
      \end{center}
      \vspace{.2cm}
    \end{minipage}

  }%
  \hfill\begin{minipage}{.20\textwidth}
    \includegraphics[width=\textwidth]{escudoInfor}
  \end{minipage}
  \vfill
  % Datos de alumno, curso y tutores ------------------
  \begin{center}%
  {%
    \noindent\LARGE
    Presentado por \@author{}\\ 
    en Universidad de Burgos --- \@date{}\\
    Tutor: \@tutor{}\\
  }%
  \end{center}%
  \null
  \cleardoublepage
  }
\makeatother


% Datos de portada
\title{Aplicación de clustering de datos educacionales en Moodle \\Documentación Técnica}
\author{Xing Long Ji}
\tutor{Dr. D. Raúl Marticorena Sánchez}
\date{\today}

\begin{document}

\maketitle



\cleardoublepage



%%%%%%%%%%%%%%%%%%%%%%%%%%%%%%%%%%%%%%%%%%%%%%%%%%%%%%%%%%%%%%%%%%%%%%%%%%%%%%%%%%%%%%%%



\frontmatter


\clearpage

% Indices
\tableofcontents

\clearpage

\listoffigures

\clearpage

\listoftables

\clearpage

\mainmatter

\appendix

\apendice{Plan de Proyecto Software}

\section{Introducción}

\section{Planificación temporal}

\section{Estudio de viabilidad}

\subsection{Viabilidad económica}

\subsection{Viabilidad legal}



\apendice{Especificación de Requisitos}

\section{Introducción}

En este anexo se comentará los objetivos generales del proyecto y detallando sobre los requisitos funcionales y no funcionales establecidos durante el desarrollo del proyecto.

\section{Objetivos generales}
El proyecto tiene como objetivos la refactorización del código junto con la integración de los registros del curso. Añadir nuevas mejoras de interacción con el usuario así como poder guardar los datos en ficheros cifrados. También realizar nuevas gráficas de los registros junto con sus opciones de filtrado.

\section{Catalogo de requisitos}

Estos son los siguientes requisitos funcionales y no funcionales:

\subsection{Requisitos funcionales}

\begin{itemize}
	

	\item \textbf{RF-1 Guardar los datos:} la aplicación tiene que ser capaz de guardar los datos para sus posteriores ejecuciones.
	\begin{itemize}
		\item \textbf{RF-1.1 Almacenar en un fichero:} la aplicación debe almacenarlo en archivos locales.
		\item \textbf{RF-1.2 Cifrar los datos:} la aplicación debe encriptar los datos cuando se guarde.
	\end{itemize}
	\item \textbf{RF-2 Recordar datos sesión:} la aplicación debe dar la opción de recordar el usuario y \textit{host} en sucesivas ejecuciones.
	\item \textbf{RF-3 Opción de actualizar:} el usuario debe tener la posibilidad de elegir 
	\item Visualización de los registros con diferentes filtros
	\item Aspectos de visualización de los usuarios
	\item Continuación con proyectos previos
	
	
	
\end{itemize}

\subsection{Requisitos no funcionales}

\begin{itemize}
	\item \textbf{RNF-1 Tiempos de respuesta:} la aplicación debe tener buenos tiempos de respuesta a las diferentes acciones del usuarios. 
	\item \textbf{RNF-2 Gestión de memoria:} el consumo de memoria de la aplicación no debe ser excesiva.
	\item Usabilidad de la aplicación.
\end{itemize}


\section{Especificación de requisitos}

\imagenflotante{diagrama_casos_uso}{Diagrama de casos de uso.}
\apendice{Especificación de diseño}

\section{Introducción}

En este anexo se va a comentar los aspectos más relevantes del diseño de esta aplicación.

\section{Diseño de datos}

Toda la información que se almacena en Moodle es con bases de datos con un campo numérico de identificador único (\textit{ID}) para cada uno de los elementos (usuarios, cursos, módulos del curso). Para simular este efecto en Java se ha creado una clase que incluye varios HashMap para los elementos comentados anteriormente usando los identificadores de cada tipo como clave del HashMap tal y como se puede ver en la figura \ref{fig:database_java}. Está clase también es el encargado de instanciar cada uno de los tipos con el objetivo de tener instancias únicas por cada identificador.

\imagen{database_java}{Simulación de base de datos.}


\section{Diseño procedimental}

En este apartado informa sobre las conexiones realizadas en la plataforma de Moodle durante la ejecución del producto. Comunicándose además de los servicios web de REST API, se inicia sesión también en la web para realizar acciones no disponibles como la descarga de los registros. Todas las peticiones HTTP y recibir las respuestas lo realiza librería jsoup.

Estas comunicaciones lo dividiremos en dos partes:

\subsection{Inicio de sesión}

Al introducir los datos de usuario, contraseña y \textit{host} en la aplicación, se inicia sesión tanto en el servicio de REST API como en la web para \textit{web scraping}.

Los pasos que sigue la aplicación en el inicio de sesión son:
\begin{enumerate}
	\item Busca el \textit{token} del REST API a partir del usuario y contraseña como parámetros de URL.
	\item Busca el \textit{login token} de inicio de sesión en la web.
	\item Inicia sesión usando el \textit{login token} como parámetro adicional si existe, guarda los \textit{cookies} de sesión para futuras peticiones. 
	\item En la API recupera la información del usuario conectado usando la función del servicio web \textbf{core\_user\_get\_users\_by\_field}.
	\item Descarga la imagen del usuario conectado.
	\item Recupera de los servicios web la información de los cursos matriculados por el usuario con la función \textbf{core\_enrol\_get\_users\_courses}.
	\item Con la función de la API \textbf{core\_course\_get\_categories} devuelve la información de las categorías en el que se encuentran los cursos.
\end{enumerate}
Estos pasos se visualizan también en el diagrama de secuencias de la figura \ref{fig:diagrama_secuencia_login}.

\imagenflotante{diagrama_secuencia_login}{Diagrama de secuencia del inicio de sesión.}



\subsection{Elección del curso}

Después de elegir el curso, la aplicación descargará información de la REST API y también los registros del curso. 

\subsubsection{REST API}

De la información que se busca en los servicios web de Moodle:


\begin{enumerate}
	\item A través de la función \textbf{core\_enrol\_get\_enrolled\_users} recogemos la información de los usuario matriculados del curso junto con sus roles y grupos.
	\item Descargamos cada uno de las fotos de los usuarios matriculados.
	\item Se busca la información de los módulos del curso que hay con la función
	\textbf{core\_course\_get\_contents}.
	\item Con la función \textbf{gradereport\_user\_get\_grades\_table} nos proporciona los calificadores de todos los alumnos.
	
\end{enumerate}

Esta información queda reflejada en el diagrama de secuencia de la figura \ref{fig:diagrama_secuencia_webservice}
\imagenflotante{diagrama_secuencia_webservice}{Diagrama de secuencia de recogida de datos de los servicios web.}

\subsubsection{Descarga de registros}

En la descarga de los registros es diferente si tiene el fichero de caché guardado previamente o no.

\begin{itemize}
	\item Cuando existe el archivo caché
	\begin{enumerate}
		\item Si la zona horaria del usuario es la misma que el servidor (99) hay que buscarlo en el perfil del usuario mediante \textit{web scraping}. 
		\item Luego va descargando los registros de forma diaria desde la fecha del último registro hasta la fecha y hora actual.
	\end{enumerate}
	\item Cuando no existe el archivo caché
	\begin{enumerate}
		\item Busca la zona horaria del servidor a partir del perfil del usuario.
		\item Descarga los registros completos del curso.
	\end{enumerate}
\end{itemize}

En el diagrama de secuencia de la figura \ref{fig:diagrama_secuencia_logs} recoge los pasos que se realiza en la descarga de registros.

\imagen{diagrama_secuencia_logs}{Diagrama de secuencia de descarga de los registros del curso.}

\section{Diseño arquitectónico}

\subsection{Arquitectura cliente-servidor}
Esta arquitectura consiste en un cliente que realiza peticiones a un servidor que le da respuestas. En este caso, el cliente sería la aplicación y el servidor es la plataforma de Moodle y sus servicios web (REST API) y la web.

En la figura \ref{fig:diagrama_despliegue} se visualiza un diagrama de despliegue que muestra la conexión de la aplicación con el servidor Moodle a través de Internet.


\imagen{diagrama_despliegue}{Diagrama de despliegue de la aplicación.}

\subsection{Arquitectura MVC}

\subsection{Patrones de diseño}

\section{Diseño de la interfaz}
En cuanto al diseño de la interfaz se ha mantenido el del proyecto de UBUGrades añadiendo las nuevas características del presente desarrollo.
\apendice{Documentación técnica de programación}

\section{Introducción}

En esta sección del anexo se explicarán cómo están organizados los paquetes del proyecto; un manual para el programador sobre las herramientas que debe instalar para seguir el desarrollo; su compilación, instalación y ejecución del proyecto y las pruebas realizadas en la aplicación.

\section{Estructura de directorios}

El repositorio del proyecto está distribuido de la siguiente forma:

\begin{itemize}
	\item /doc/: los ficheros javadoc del proyecto.
	\item /latex/: los elementos necesarios para generar la memoria y el anexo, tanto los ficheros .tex como los .pdf. También se incluye las imágenes y la bibliografía en formato .bib.
	\item /lib/: librerías externas de Java.
	\item /lib/src: el código fuente de las librerías externas de Java.
	\item /lib/docs: los javadoc de las librerías externas de Java.
	\item /python/: varios archivos de utilidad para los registros del curso:
	\begin{itemize}
		\item \textbf{Componentes y eventos.json}: fichero dividido en tres niveles que son Componentes, Eventos y Descripciones de los registros del curso.
		\item \textbf{Contains checker.ibynb}: funciones de utilidad que comprueba en las Descripciones de los registros si existe o no un carácter o varios.
		\item \textbf{Generador de código componentes y eventos.ipynb}: funciones que generan código Java automáticamente a partir del fichero JSON mencionado.
		\item \textbf{Rastreador de componentes y eventos.ipynb}: programa de utilidad que escanea los registros del curso en formato CSV y con eso crea un fichero JSON con todos los Componentes, Eventos y Descripciones encontrados.
	\end{itemize}

	\item /resources/: recursos de apoyo para el programa.
	\item /resources/css/: ficheros CSS que modifica la visualización y colores de la aplicación.
	\item /resources/graphics/: ficheros HTML para dibujar las gráficas.
	\item /resources/graphics/lib: librerías externas de JavaScript usado para las gráficas.
	\item /resources/img/: todas las imágenes de la aplicación.
	\item /resources/messages/: ficheros de internacionalización de textos.
	\item /resources/view/: se encuentra los ficheros FXML.
	\item /src/: código fuente del desarrollo.
	\item /src/controllers/: paquete encargado del flujo de la ejecución.
	\item /src/controllers/datasets/: dirige las generaciones de los datos en formato válido para Chart.js.
	\item /src/controllers/ubugrades/: encargado de la generación de los cursos, usuarios matriculados, calificaciones, etc. Toda la parte de los servicios web.
	\item /src/controllers/ubulogs/: gestiona la descarga de los registros y su posterior \textit{parseo}.
	\item /src/conrollers/ubulogs/logtypes/: las clases encargadas de gestionar los ids de las Descripciones del registro.
	\item /src/model/: representación de los datos, su lógica de negocio.
	\item /src/model/mod: todas las clases de los tipos de módulos del curso.
	\item /src/persistence/: encriptación y desencriptación de archivos de objetos Serializados.
	\item /src/webservice/.../: paquete y subpaquetes encargados de llamar a las funciones de Moodle añadiendo los parámetros necesarios a la URL y recibir las respuestas.
	

\end{itemize}

\section{Manual del programador} \label{sec:manual_programdor}

Este manual sirve de referencia para personas que tomen el proyecto en el futuro. Se explicarán como construir el entorno de desarrollo, qué es lo que necesitan instalar.

\subsection{Instalación de Java}

La aplicación está desarrollada en Java 8, por lo tanto para trabajar en ella necesitará tener el JDK 8. Durante el proyecto se ha trabajado con la \textbf{versión 1.8.0\_201} pero sería recomendable trabajar con la última disponible. Realizamos los siguientes pasos si queremos descargar la última versión de JDK: 

\begin{enumerate}
	\item Entrar en la página de  \href{https://www.oracle.com/technetwork/java/javase/downloads/index.html}{descargas de Java}.
	\item Buscamos \textbf{Java SE 8uxxx} siendo las xxx el número de actualización de Java 8. 
	\item Descargamos el JDK.
	\imagen{elegir_jdl}{Elegir la descarga de JDK.}
	\item Leemos la licencia de uso de Oracle y aceptamos las condiciones.
	\item Cuando hayamos aceptado la licencia de uso descargamos el ejecutable según el sistema operativo que tengamos.
	
	\textbf{Nota:} x86 es de 32 bits y x64 de 64 bits.
	\imagen{descarga_java}{Descarga del JDK 8}
	\item Seguir las instrucciones del instalador
	
	
\end{enumerate}

Es recomendable desinstalar la versión antigua de Java si descargamos la última versión \cite{noauthor_por_nodate}.
\subsection{Eclipse}

Durante el desarrollo se ha empleado la versión \textbf{Oxygen.1a Release (4.7.1a)} de Eclipse aunque no habría muchos problemas usar una versión más reciente. A continuación se muestra el método de descargar una versión portable de Eclipse:

\begin{enumerate}
	\item  Entrar en la página web descarga de \href{https://www.eclipse.org/downloads/packages/}{paquetes de Eclipse}.
	\item Descargamos el paquete \textbf{Eclipse IDE for Java Developers}, en el sistema operativo que se tenga.
	\imagen{descargar_eclipse}{Descarga del paquete de Eclipse.}
	\item Pulsamos el botón de descarga \textbf{Download}.
	\imagen{boton_descarga_eclipse}{Pulsar el botón de descarga.}
	\item Descomprimir la carpeta dentro del ZIP.
	\item Ejecutar \textbf{eclipse.exe} cuando se quiera usar.
\end{enumerate}


\subsection{Instalación de e(fx)clipse}

e(fx)clipse es una herramienta adicional de Eclipse que facilita el desarrollo de aplicaciones que usen JavaFX. 
Para instalar e(fx)clipse realizamos los siguientes pasos:

\begin{enumerate}
	\item Seleccionamos el menú \textbf{Help} de la barra de herramientas de Eclipse.
	\item Seleccionamos el elemento del menú \textbf{Install New Software...}
	
	\imagen{efxclipse}{Menú \textbf{Help} y seleccionar \textbf{Install New Software...}}
	
	\item Seleccionamos la opción \textbf{--All Available Sites --} de la lista desplegable de \textbf{Work with:}.
	\item Escribimos e(fx)clipse en el cuadro de texto.
	\item Marcamos solo la casilla que pone \textbf{e(fx)clipse - IDE}.
	\item Pulsamos el botón \textbf{Next} y debería descargar e instalarse.
	\imagen{efxclipse2}{Instalación de la herramienta.}
\end{enumerate}


\subsection{Instalación de Scene Builder}

La herramienta del Scene Builder (2.0) se puede conseguir desde: 

\href{http://www.oracle.com/technetwork/java/javase/downloads/javafxscenebuilder-1x-archive-2199384.html}{http://www.oracle.com/technetwork/java/javase/downloads/javafxscenebuilder-1x-archive-2199384.html}

Después de entrar a la página se busca la versión 2.0 de Scene Builder, se acepta la licencia y descargamos la versión según el sistema operativo que se use.

\imagen{descargar_scene_builder}{Aceptar la licencia de Oracle BSD y descargr.}

Se sigue las instrucciones del instalador.


\subsection{Instalación de Moodle}

Para instalar Moodle en el equipo personal, se puede descargar desde la página oficial con todo incluido para su funcionamiento:

\href{https://download.moodle.org/windows/}{https://download.moodle.org/windows/}

Una vez descargado y descomprimido la carpeta, ejecutamos \textbf{StartMoodle.exe} que inicia la  arranque del servidor.

Para acceder al sitio web de Moodle se usa el localhost: \href{http://localhost}{http://localhost}
En la página web, la primera vez iniciará el proceso de instalación del servidor. Se sigue las instrucciones del instalador.

Por defecto no están activados los servicios web de Moodle. Para habilitarlo seguimos los siguientes pasos después de iniciar sesión con la cuenta del administrador:

\begin{enumerate}
	
	\item Pulsamos el botón \textbf{Administración del sitio}.
	\item Entramos dentro de \textbf{Mobile Settings}.
	\imagen{servicios_moodle}{Habilitar servicios web.}
	\item Marcamos la casilla de \textbf{Habilitar servicio web para dispositivos móviles}.
	\item Pulsamos en \textbf{Guardar cambios}
	\imagen{servicios_moodle2}{Habilitar servicios web parte 2.}
\end{enumerate}

\subsubsection{Listado de funciones de los servicios web de Moodle}

Moodle ofrece también un listado de funciones del servicio web junto con información de parámetros y el esqueleto del JSON que devuelve. Para acceder a ella es necesario haber iniciado sesión con la cuenta del administrador.
\begin{enumerate}
	\item Después de iniciar sesión con el administrador. Pulsamos \textbf{Administración del sitio} y luego a la pestaña de \textbf{Extensiones}.
	\imagen{api_moodle}{Acceso al listado de funciones del servicio web de Moodle.}
	\item Bajamos abajo hasta la fila de \textbf{Servicios Web} y accedemos a \textbf{Documentación API}-
	\imagen{documentacion_api_moodle}{Acceder a la Documentación API.}
	\item Muestra un listado de todas las funciones que al desplegar se puede visualizar los parámetros y esqueletos de la información que devuelve.
	\imagen{listado_de_funciones}{Listado de funciones del servicio web de Moodle.}
\end{enumerate}

\subsection{Instalación de Jupyter Notebook}

Según la página oficial de Jupyter Notebook \cite{noauthor_project_nodate}, lo más recomendable es instalarlo a través de Anaconda en:

\href{https://www.anaconda.com/distribution/\#download-section}{https://www.anaconda.com/distribution/\#download-section}.

La versión recomendable que se debe descargar es el que incluye Python 3.7 o superior debido a que los diccionarios están ordenados por orden de inserción.

\subsection{Clasificación de registros}

Para clasificar se ha partido de la premisa de que las columnas \textbf{Componente} y \textbf{Nombre de Evento} usa un esqueleto de \textbf{Descripción} único. 

A partir de expresiones regulares, se ha buscado las posibles los ids de las descripciones y con esos números pasado a un delegado (clase) en función del Componente y Evento. 

\imagen{tipos_logs}{Clasificación de los registros junto con sus delegados.}

\subsubsection{Delegados}

Los nombres de los delegados son códigos nemotécnicos en función el tipo de ids asigna. Por ejemplo en la siguiente imagen usa los códigos nemotécnicos User, Affected y Cmid. Está asignando el usuario, usuario afectado y el módulo del curso.
\imagen{ejemplo_tipo_log}{Un delegado UserAffectedCmid}

En la tabla \ref{tabla:tipos_delegados} se puede ver los códigos nemotécnicos usados y una pequeña descripción.

\tablaSmall{Códigos nemotécnicos de los delegados.}{l c }{tipos_delegados}{ \multicolumn{1}{l}{Código nemotécnico} & Descripción \\}{ 
	Affected         & id del usuario afectado                          \\
	Attempt          & id del intento                                   \\
	Calendar         & id del evento del calendario                     \\
	Category         & id de la categoría                               \\
	Chapter          & id del capitulo del libro                        \\
	Choice           & id de elección (choice)                          \\
	Cmid             & id del módulo del curso                          \\
	Comment          & id del comentario                                \\
	Competency       & id de competencia del curso                      \\
	Course           & id del curso                                     \\
	Discussion       & id de la discusión del foro                      \\
	Evidence         & id de la evidencia                               \\
	Field            & id del campo                                     \\
	Files            & número de ficheros                               \\
	Glossary         & id de la entrada al glosario                     \\
	Grade            & id de calificación                               \\
	Gradeitem        & id de grade ítem                                 \\
	Group            & id del grupo                                     \\
	Grouping         & id de grouping                                   \\
	Item             & id de item                                       \\
	Note             & id de note                                       \\
	Option           & id de opción                                     \\
	Override         & id de anulación (override)                       \\
	Page             & id de la página                                  \\
	Post             & id del post del foro                             \\
	Question         & id de pregunta del cuestionario                  \\
	Questioncategory & id de categoría de cuestionario                  \\
	Record           & id de record                                     \\
	Role             & id del rol                                       \\
	Rule             & id de regla                                      \\
	Scale            & id de escala                                     \\
	Sco              & id del scorm package                             \\
	Section          & número de sección                                \\
	Submission       & id de la entrega                                 \\
	Subscription     & id de subscripción al evento                 	\\
	Tour             & id de tour                                       \\
	User             & id del usuario del usuario que realiza la acción \\
	Words            & número de palabras                               \\
}

\subsection{Uso de las herramientas de utilidad Python}

En este apartado se comentará las tres herramientas existentes en Jupyter Notebook. Estos facilitan al programador crear ciertos códigos que tienen que ver con los registros del curso.

\subsubsection{Rastreador de componentes y eventos} \label{rastreador_componentes_eventos}

Programa que selecciona uno o varios CSV de los registros descargados en inglés. Y de forma opcional elegir el archivo JSON donde se actualizarán todos los \textbf{Componentes}, \textbf{Eventos} junto con las \textbf{Descripciones}. Si no se selecciona el fichero JSON generará uno nuevo de los registros del CSV.

Se aporta ya un fichero JSON (\textbf{Componentes y eventos.json}) con \textbf{48 Componentes} y \textbf{193 Eventos} diferentes detectados hasta el momento. 

Este fichero usa tres niveles:
\begin{enumerate}
	\item En el primer nivel indica el Componente.
	\item En el segundo nivel es el Evento.
	\item En el tercer nivel es un listado del esqueleto de las Descripciones asociados al Componente y Evento.
\end{enumerate}

En la figura \ref{fig:json_componente_evento} podemos ver un extracto del JSON recuadrado los Componentes (en rojo) y los Eventos (en verde).

\imagen{json_componente_evento}{Extracto de Componentes y eventos.json}

Los pasos a seguir añadir nuevos elementos son:

\begin{enumerate}
	\item Abrir el archivo \textbf{Rastreador de componentes y eventos.ipynb} con Jupyter Notebook.
	\item El programa está preparado para ejecutar todas las celdas seguidas, elegimos el menú \textbf{Cell} y luego la opción \textbf{Run All}.
	\imagen{rastreador}{Menú \textbf{Cell} y opción \textbf{Run All}}
	\item Aparecerá una primera ventana que pide seleccionar los ficheros CSV de los registros en inglés.
	\imagen{elegir_Csv}{Selección múltiple de los archivos CSV.}
	\item A continuación se mostrará una nueva ventana que pide el fichero JSON.
	\imagen{seleccion_json}{Seleccionar el archivo JSON.}
	\item Si no hay ningún problema encontrado durante la ejecución de las celdas, ha salido todo bien. 
\end{enumerate}

Uno de los problemas que se puede encontrar es que haya elegido accidentalmente un CSV en otro idioma que no sea en inglés. Una de las celdas saltará un error y las siguientes no se ejecutarán.

Por ejemplo en la siguiente imagen se puede ver que salta un error de que no se encuentra la columna Component. Producido por el último fichero en \textit{parsearse} indicado en un recuadro en rojo.
\imagen{key_error}{Error producido al usar un CSV que no está en inglés.}


Para solucionar solo hay que volver ejecutar todas las celdas sin elegir ese archivo CSV que da problemas.

\subsubsection{Generador de código componentes y eventos}

Este segundo programa usa el fichero JSON generado comentado anteriormente (\ref{rastreador_componentes_eventos}) y va mostrando en pantalla lo que hay que copiar y pegar en el proyecto.

Genera automáticamente las enumeraciones de los Componentes y Eventos, las traducciones del \textit{Resource Bundle} en inglés y las combinaciones de Componente y Evento existentes.

\begin{enumerate}
	\item Abrir el archivo \textbf{Generador de código componentes y eventos.ipynb} con Jupyter Notebook.
	\item El programa está preparado para ejecutar todas las celdas seguidas, elegimos el menú \textbf{Cell} y luego la opción \textbf{Run All}.
	\imagen{generador_java}{Menú \textbf{Cell} y opción \textbf{Run All}}
	\item Lo primero que mostrará una ventana que pide seleccionar el archivo JSON.
	\imagen{seleccion_json}{Seleccionar el archivo JSON.}
	\item En las siguientes celdas se parará en una que pide que copiemos los componentes del Resource Bundle en español desde:
	
	\textbf{/resources/messages/messages\_es.properties}
	\imagen{resource_bundle_es}{Celda que parará el flujo el ejecución.}
	\imagen{copiarcomponentes}{Copiar los componentes que se encuentra en el Resource Bundle español.}
	\item Cuando lo hayamos copiado, en el cuadro de texto Notebook pulsamos la tecla \textbf{Entar} para que siga ejecutando las siguientes celdas.
	\item Realizamos las mismas acciones pero ahora con los eventos del Resource Bundle en español.
	\imagen{copiar_eventos}{Copiar los eventos que se encuentra en el Resource Bundle español.}
	\item Realizamos la misma acción con los tipos de logs que se encuentra en:
	
	 \textbf{/controllers/ubulogs/logtypes/Logtypes.java}
	\imagen{copiar_logtypes}{Copiar los tipos de logs.}
	\item El último paso es solo ir copiando y pegando las salidas de las celdas al proyecto de Java.
\end{enumerate}



\section{Compilación, instalación y ejecución del proyecto}

\subsection{Importar el proyecto a Eclipse}
Lo primero de todo hay que descargar todo el ficheros fuentes desde el repositorio de GitHub y descargamos el ZIP:

\href{https://github.com/yjx0003/UBUMonitor}{https://github.com/yjx0003/UBUMonitor}

Para importar en Eclipse seguimos los siguientes pasos:

\begin{enumerate}
	\item Pulsar el menú \textbf{File} y elegir la opción \textbf{Import...}
	\imagen{import_project}{Menú File y opción Import...}
	\item Seleccionar la opción de \textbf{Existing Projects into
		Workspace} y luego al botón \textbf{Next}.
	\imagen{existing_project}{Existing Projects into
		Workspace} 
	\item Marcamos la opción \textbf{Select archive file} y seleccionamos el fichero ZIP descargado previamente del repositorio GitHub. Pulsamos Finish e importara el fichero.
	
	\imagen{import_projects}{Selección del archivo ZIP e importar proyecto.}
\end{enumerate}

\subsection{Generación del ejecutable de la aplicación}

Para generar el ejecutable jar de la aplicación debemos seleccionar la opción \textbf{Exportar...} con click derecho del ratón al proyecto.
\imagen{export_jar}{Exportar jar}

Seleccionamos la exportación del proyecto como \textbf{Runnable JAR file} dentro de la carpeta Java.

\imagen{runnable_jar}{Seleccionar Runnable JAR file.}

Seleccionar al opción \textbf{Extract required libraries into generated JAR} y exportar.

\imagen{runnable_jar_2}{Configuración del Runnable JAR.}

\subsection{Ejecución de la aplicación}

Para ejecutar el proyecto desde Eclipse: Buscamos el archivo java \textbf{UBUMonitor} dentro del paquete \textbf{controllers} que es donde se encuentra el Main. Click derecho y \textbf{Run as...} -> \textbf{Java Application}

\imagen{run}{Ejecución de la aplicación.}

La ejecución del jar ejecutable es con doble \textit{click} directamente.
\section{Pruebas del sistema}

Las pruebas de las nuevas funcionalidades se realizaron contra la el Moodle instalado en el equipo en la versión \textbf{3.5.4+} y el servidor \textit{demo} de Moodle \href{https://school.demo.moodle.net}{Mount Orange School} en la versión \textbf{3.7}. 

Cuando se tenía unos cambios grandes, los tutores realizaban pruebas contra UBUVirtual con datos reales de alumnos. Con esto daban el visto bueno o informes de errores encontrados.
\apendice{Documentación de usuario}

\section{Introducción}
En este apartado se indicarán los requisitos necesarios para utilizar la aplicación, ademas de una guía de instalación. Se incluye un manual de usuario para el uso de la aplicación.

\section{Requisitos de usuarios}
\begin{itemize}
	\item Tener el \textbf{JRE de Java 8} instalado. Se puede utilizar el de \href{https://www.java.com/es/download/}{Oracle} o el de \href{https://www.azul.com/downloads/zulu-community/?version=java-8-lts&os=windows&architecture=x86-64-bit&package=jre}{Zulu}. Se recomienda utilizar la distribución de Zulu.
	\item Conexión a internet.
	\item Acceso a un servidor Moodle con rol de profesor.
\end{itemize}
\section{Instalación}
Para la instalación del programa solo hay que descargarse el \textbf{jar} desde el \href{https://github.com/xjx1001/UBUMonitor/releases/latest}{repositorio de GitHub}.
Para ejecutar el \textbf{jar} hay que hacer doble click en el archivo, o mediante la consola del sistema. Para ejecutar mediante la consola:
\begin{enumerate}
	\item Abrimos la ubicacion del archivo jar en el explorador de archivos.
	\item Escribimos \textbf{cmd} en la ruta y pulsamos \textbf{Enter} .
	\imagen{anexos/instalacion1}{Escribir \textbf{cmd} en el recuadro marcado.}
	\item Introducir el comando: \texttt{java -jar UBUMonitor-Clustering-vX.X.jar}
	\item Presionar \textbf{Enter}.
\end{enumerate}

\section{Manual del usuario}
En este manual se explicaran todas las funcionalidades que tiene la aplicación.

Para iniciar el programa ejecutamos el \emph{jar} que habíamos descargado.

\subsection{Inicio de sesión}
Al ejecutar la aplicación, la primera pantalla que se muestra es la del inicio de sesión. Para esto se pide el usuario, la contraseña y la url del servidor Moodle.

\imagenConTamano{anexos/manual1}{Inicio de sesión.}{0.6}

En esta pantalla también se puede modificar el idioma de la aplicación, en el menú desplegable de arriba. La opción de \textbf{modo offline} no puede ser marcada en la primera ejecución, ya que recoge los datos almacenados de forma local.

Una vez rellenado los datos pulsamos en el botón \textbf{Entrar}.

\subsection{Elección del curso}
En esta pantalla se muestra los cursos a los que tiene acceso el usuario. Si estamos en \textbf{modo \emph{offline}} solo se mostrarán aquellos cursos que se han accedido previamente. Si estamos en \textbf{modo \emph{online}} nos permitirá actualizar los datos de un curso, esta actualización es obligada si es la primera vez que se accede al curso.

\imagenConTamano{anexos/manual2}{Selección de un curso en modo online.}{0.7}

\subsection{Clustering}
A partir de aquí solo se explicará las funcionalidades refrentes a este proyecto (la parte de \emph{clustering}). Para ello, una vez accedido a un curso, presionamos en la pestaña de \textbf{\emph{clustering}}.
\imagen{anexos/manual3}{Pestaña de \emph{Clustering}.}

\subsubsection{Ejecución del clustering} \label{sssection:ejecutar}
Para comenzar la ejecución, primero hay que seleccionar los alumnos con los que se quiere realizar el \emph{clustering}. Para seleccionar los alumnos hacemos \texttt{Ctrl + click izquierdo} sobre el alumno. Es posible realizar filtros de usuarios en función del rol, grupo o último acceso al curso.
\imagenConTamano{anexos/manual4}{Selección de alumnos.}{0.5}

Después, hay que elegir con qué datos proporcionados por Moodle queremos trabajar. Para ello seleccionamos de igual manera que en la selección de alumnos. Se pueden seleccionar elementos de varias pestañas.
\imagenConTamano{anexos/manual5}{Selección de datos.}{0.5}

A continuación, hay que seleccionar el algoritmo y sus parámetros de ejecución. Cada algoritmo tiene diferentes parámetros, por lo que primero hay que seleccionar el algoritmo y luego los parámetros.
\imagen{anexos/manual6}{Selección del algoritmo y sus parámetros.}
El \textbf{\emph{tooltip}} de cada parámetro da una pequeña explicación. 

Ahora, hay que marcar los tipos de datos con los que se van a realizar el \emph{clustering}, marcar un tipo de dato significa que se van a utilizar los datos seleccionados anteriormente. Si se seleccionan los registros, hay que marcar los tipos, que por defecto están marcados todos, y en intervalo de días. El inicio es el selector de fechas superior y el final es el inferior, por defecto esta el inicio de curso configurado en Moodle y el día actual. Esta opción es útil para realizar el \emph{clustering} en un instante de pasado concreto.
\imagenConTamano{anexos/manual7}{Tipos de datos.}{0.8}

Además, se puede marcar la opción de filtrar datos, el cual elimina los valores constantes a todos los alumnos seleccionados antes de realizar el \emph{clustering}. Es recomendable marcar esta opción.

Para finalizar, hay que decidir cuántas iteraciones del algoritmo se van a ejecutar. También existe la posibilidad de realizar una reducción de dimensionalidad, esta opción reduce la cantidad de datos, manteniendo solo los valores mas significativos. Puede ser útil si se trabajan con una gran cantidad de datos, esta opción reduciría el tiempo de ejecución.
\imagenConTamano{anexos/manual8}{Tipos de datos.}{0.6}

\subsubsection{Errores al ejecutar}
A continuación se van a enumerar los posibles errores al ejecutar:
\begin{itemize}
	\item No hay usuarios seleccionados: se produce cuando no se ha seleccionado ningún usuario o solo un usuario.
	\imagenConTamano{anexos/manual9}{No hay usuarios seleccionados}{0.6}
	\item No hay datos seleccionados: se produce cuando no se ha seleccionado ningún elemento de las listas marcadas en los tipos de datos. También se produce cuando no se ha marcado ningún tipo de dato.
	\imagenConTamano{anexos/manual10}{No hay datos seleccionados}{0.6}
	\item Error en un parámetro del algoritmo: se produce cuando se ha introducido un valor inválido en algún parámetro.
	\imagenConTamano{anexos/manual11}{Error en un parámetro del algoritmo }{0.6}
\end{itemize}

\subsubsection{Resultado del clustering}
Cuando termine de ejecutar el resultado de la ejecución se mostrará en gráficas y en tabla. Inicialmente los nombres de los clústeres son números.

La gráfica 2D es un diagrama de dispersión donde cada punto es un alumno y el color representa el clúster al que pertenece, adicionalmente se muestra el centroide de cada clúster si el algoritmo ejecutado no está basado en densidades.
\imagen{anexos/manual12}{Diagrama de dispersión 2D.}

La gráfica 3D es similar a la 2D, salvo que no se muestran los centroides.
\imagen{anexos/manual13}{Diagrama de dispersión 2D.}

Ambos gráficos son interactivos, se puede ocultar los puntos de un clúster presionando sobre la leyenda.

En el análisis de silueta se muestra el diagrama de barras representado el análisis realizado. El valor del ancho de silueta indica lo bien que está clasificado un alumno. Cuanto mayor es el valor, mejor está clasificado. La linea gris representa la media.
\imagen{anexos/manual14}{Análisis de silueta.}

En la tabla \ref{fig:anexos/manual15} se muestran los alumnos con los que se han realizado el \emph{clustering}. Puede que no salgan todos, ya que algunos algoritmos basados en densidades eliminan el ruido. Se muestra la foto del alumno, su nombre completo y el clúster al que pertenece.
\imagenConTamano{anexos/manual15}{Tabla del \emph{clustering}.}{0.7}

En el menú desplegable se pueden filtrar los alumnos por clúster. Adicionalmente se mostrará en la tabla las calificaciones seleccionadas, las notas están en el intervalo [0, 100]. El intervalo [0, 25] es de color rojo, [25, 50] de color amarillo, [50, 75] de color verde y [75, 100] de color morado.
\imagenConTamano{anexos/manual16}{Tabla con calificaciones.}{0.7}

En la tabla al hacer doble-click o click derecho sobre un alumno nos mostrará información de este alumno utilizado en el \emph{clustering}.
\imagen{anexos/manual19}{Información de un alumno.}

Las flechas de abajo sirven para pasar al siguiente o anterior alumno. En la tabla se marca al alumno que se está visualizando actualmente.

\subsubsection{Modificar los nombres}
La aplicación permite modificar los nombres de los clústeres, esta modificación será efectiva en todas las gráficas y en la tabla. Para realizar esto vamos a la pestaña \textbf{Etiquetar}.
\imagenConTamano{anexos/manual17}{Etiquetar.}{0.5}

Las etiquetas que utilicemos se guardarán, y aparecerán en el autocompletar cuando las volvamos a escribir.

La etiquetas creadas se pueden eliminar en la pestaña \textbf{Gestionar etiquetas}.
\imagenConTamano{anexos/manual18}{Gestionar etiquetas.}{0.5}

\subsubsection{Realizar un análisis}
Para ejecutar un análisis hay que seguir los mismos pasos que para ejecutar un algoritmo. El análisis solo se puede realizar con algoritmos que tengan el paramento \textbf{Nº de agrupaciones}. Solo cambia el ultimo paso, en vez de presionar al botón de ejecutar hay que presionar el de analizar. El selector de rangos indica con cuántas agrupaciones se tiene que realizar el \emph{clustering}. 

Existen dos tipos de análisis, el del codo y el de silueta.

\imagenConTamano{anexos/manual20}{Análisis del codo.}{0.6}
\imagenConTamano{anexos/manual21}{Análisis de silueta.}{0.6}

\subsection{Clustering jerárquico}
Para realizar una agrupación jerárquica hay que cambiar de pestaña, de particional a jerárquico.
\imagen{anexos/manual22}{Pestaña del \emph{clustering} jerárquico.}

Antes de ejecutar hay que seleccionar los alumnos y los datos al igual que en el \emph{clustering} particional \ref{sssection:ejecutar}.

Se puede modificar la medida de distancia entre elementos y también la distancia entre agrupaciones.
\imagenConTamano{anexos/manual23}{Selección de las medidas de distancia.}{0.7}

Después de ejecutar se mostrara el dendrograma resultante de la ejecución.
\imagen{anexos/manual24}{Dendrograma resultante.}

Una vez ejecutado el \emph{clustering} jerárquico se pueden realizar particiones en cualquier número de agrupaciones.
\imagenConTamano{anexos/manual25}{Selección del número de agrupaciones}{0.7}

Una vez ejecutado se mostrara el resultado en la tabla. Esta tabla es la misma que en el \emph{clustering} particional, tiene las mismas funciones.
\imagenConTamano{anexos/manual26}{Tabla del resultado.}{0.7}

\subsection{Exportación de los resultados}
Todas la gráficas se pueden exportar a CSV y a PNG, la tabla se puede exportar solo a CSV. Para exportar cualquier gráfica hay que hacer click derecho en la gráfica y seleccionar el tipo de exportación. Después hay que indicar dónde queremos guardar el fichero de la exportación.

Para la tabla, se pueden exportar los datos de calificación mostrados en esta, para ello hay que marcar la opción \textbf{Exportar calificaciones}. Esta opción esta deshabilitada si en la tabla no se muestra ninguna calificación.


\bibliography{bibliografiaAnexos}
\bibliographystyle{plainnat}

\end{document}
